\chapter{Testing and Debugging}

\section{Why Test?}\index{testing}\index{debugging}
Programs fail in small ways. A function that returns the wrong answer for one
input can look fine for every other input you try by hand. Testing means writing
code that checks your code, so you find errors before someone else does.
Debugging means finding the cause of an error once you know it exists.

\section{Print Debugging}\index{debugging!print statements}
The quickest way to find a bug is to print the values that matter. Suppose
\pyfunction{average} returns the wrong answer:

\begin{minted}{python}
def average(values):
    total = sum(values)
    print("total:", total, "count:", len(values))
    return total / len(values)

print(average([1, 2]))
\end{minted}

The printed values show where the output diverges from what you expected.
Remove the print statements when you are done.

\section{Assertions}\index{assertion}
An \newterm{assertion} checks a condition and stops the program if it is false:

\begin{minted}{python}
from mathtools import average

x = [1, 2, 3]
assert average(x) == 2, "average should be 2"
print("ok")
\end{minted}

Assertions are good for stating what must always be true. They are not a
replacement for tests, because they run only when the program runs.

\section{Unit Tests}\index{unit test}
A \newterm{unit test} checks one function on known inputs and expected outputs.
Python's \texttt{unittest} module runs them for you. Given this file,
\texttt{mathtools.py}:

\begin{minted}{python}
def average(values):
    """Return the mean of a non-empty list of numbers."""
    if len(values) == 0:
        raise ValueError("average() needs at least one value")
    return sum(values) / len(values)

def is_even(n):
    return n % 2 == 0
\end{minted}

Put the tests in \texttt{test\_mathtools.py}:

\begin{minted}{python}
import unittest
from mathtools import average, is_even

class TestAverage(unittest.TestCase):
    def test_simple_list(self):
        self.assertEqual(average([2, 4, 6]), 4)

    def test_single_value(self):
        self.assertEqual(average([7]), 7)

    def test_empty_list_raises(self):
        with self.assertRaises(ValueError):
            average([])

class TestIsEven(unittest.TestCase):
    def test_even_and_odd(self):
        self.assertTrue(is_even(10))
        self.assertFalse(is_even(7))

if __name__ == '__main__':
    unittest.main()
\end{minted}

Run \texttt{python3 -m unittest test\_mathtools.py -v}. Each test prints a
line, and the final line says \texttt{OK} when every test passes.

\section{Reading a Failing Test}
A test that fails shows what it expected and what it got. If you change the
expected value in the test to 2 for \texttt{average([1, 2])}, the output
reads:

\begin{minted}{text}
AssertionError: 1.5 != 2
\end{minted}

The function is correct and the test is wrong. Before you change the code,
decide which one is wrong. Check the function's definition, then the expected
value.

\section{A Debugging Routine}
\begin{enumerate}
\item Reproduce the error with a specific input.
\item Write a test that fails for that input.
\item Form a guess about the cause, and check it with a print statement or an assertion.
\item Fix the code, and run all the tests again, not just the new one.
\end{enumerate}

\begin{Exercise}[label=test1, title=Write a test]
Add a test to \texttt{test\_mathtools.py} that checks that \texttt{average([-1, 1])} equals 0.
\end{Exercise}
\begin{Answer}[ref=test1]
\begin{minted}{python}
    def test_negative_values(self):
        self.assertEqual(average([-1, 1]), 0)
\end{minted}
This test passes, because $(-1 + 1) / 2 = 0$.
\end{Answer}

\begin{Exercise}[label=test2, title=Which test fails?]
A student wrote \texttt{self.assertEqual(average([1, 2, 3]), 3)}. What does the
test output say, and what is the mistake?
\end{Exercise}
\begin{Answer}[ref=test2]
\texttt{average([1, 2, 3])} is $\frac{6}{3} = 2$, not 3. The output says
\texttt{AssertionError: 2.0 != 3}. The expected value in the test is wrong.
\end{Answer}

\begin{Exercise}[label=test3, title=Edge case]
What should \texttt{average([])} do, and which test checks it?
\end{Exercise}
\begin{Answer}[ref=test3]
It should raise a \texttt{ValueError}. The test \texttt{test\_empty\_list\_raises}
checks this.
\end{Answer}
